\chapter{Studying cities}
\label{chap:studying_cities}

\begin{flushright}{\slshape    
Chaos was the law of nature;\\
Order was the dream of man.} \\ \medskip
--- Henry Adams~\cite{Adams:1990}
\end{flushright}

\bigskip

Cities appeared some $10,000$ years ago~\cite{Bairoch:1985, Mumford:1961}
concomitantly with the Neotlithic agriculture revolution, and really started to
thrive after the industrial revolution~\cite{Bairoch:1985}.  In England first,
where the revolution was born, with London being the first in the modern world
to reach $1,000,000$ inhabitants at the beginning of the $19$th century. The
urban growth then slowly spread through the end of the $19$th and the $20$th to
the rest of the Western world. Now, while most western countries are already
mostly urban (as of $2014$, the United States' population was $82\%$ urban,
Japan $93\%$, and most countries in the European Union were around the $80\%$
mark), most of what has been dubbed the 'urban \graffito{Source:\\ UN Population
Division (2011)} revolution' is happening in developing countries. A symbolic
barrier was reached in $2005$, when it was estimated by the U.N. that more than
$50\%$ of the world total population was living in cities. It is not difficult
to convince oneself that urbanisation is not an accident in human history, and
that cities' influence and impact are not going to stop growing any time soon.

In fact, the impact of cities is already tremendous. First, they have a
disproportionately large importance in the world's economy. A 2012 report by
McKinsey noted that while cities represented respectively $79\%$ and $19\%$ of the Unites
States' and India's population, their share in the countries' GDP was
respectively $85\%$ and $39\%$. 

Data from the NASA indicate that urban areas cover a total of $5\%$ of the total
land surface area in the world, roughly the equivalent of the superficy of the
European Union. Yet, despite their little spatial fooprint, cities have a great
impact on the environment. The United Nations indeed estimated in $2011$ that cities were
responsible for $70\%$ percent of the world's CO\textsubscript{2} emissions.

We could multiply the statistics, but the few examples given above should
convince the reader of the importance to understand cities if we want to
change the world we live in for the best. The dramatic growth of urban areas in
developing countries brings unprecendented challenges. The cause, and the
solution of some of the world's most pressing challenges find their origin in
cities. By improving the way cities work, we can hopefully make dramatic changes
to the way people live. To be able to do so however, we need to understand first
how they work.\\


\section{Cities as complex systems}
\label{sec:cities_as_complex_systems}

Cities are systems, a set of patterns and behaviours that make them identifiable
as such~\cite{Dennett:1991}. They are, in a sense, paradigmatic examples of
complex systems~\cite{Ladyman:2013}. 
First, they comprise thousands, millions of individuals that are moving and
interacting constantly. Cities are indeed more than the mere agglomeration of
residences, factories and shops in the same region; they exist and thrive
through the resulting facilitated interaction between
individuals~\cite{Bettencourt:2013,Sim:2015}. 

Second, cities are incredibly resilient systems. There are multiple examples in
History of cities that were completely destroyed -- Dresden and Hiroshima, for
instance, completely burnt to ashes during WWII -- but were later rebuilt and
thrived again.

Finally, as systems, cities exhibit very particular shapes and behaviours. The
boundaries of urban areas possess a fractal structure, we can identify them on a
map thanks to the local properties of their road network, etc. In this thesis,
we will be interested in some of the particular behaviours exhibited by cities.\\


\subsection{An organised complexity}
\label{sub:an_organised_complexity}

The systems studied in Physics can be roughly divided in two
categories~\cite{Parisi:1999}

\begin{itemize}
    \item Simple systems with only a few variables. Their dynamics is described
        by \emph{deterministic} equations. For instance, the motion of planets
        can be described with high accuracy by General Relativity.
    \item Weakly, locally interacting systems, with a very large number of
        particules. Their properties are described using \emph{probabilitistic}
        language. For instance, monoatomic gases in usual conditions of pressure
        and temperature are well described by Statistical Mechanics.
\end{itemize}

Organised randomness. 

Jane Jacosb p 489: the order is complex systems exists once we have studied it.
Order emerges from chaos because of understanding. We cannot impose order in
cities in arbitrary way, we need to understand how they work first. We can
visualise cities as a very large systems with many retroaction loops. Planning
needs to understand what these loops are, and be careful of not destroying any
of them if it wants to be succesfull. Rather than imposing a structure on cities
and the way its inhabitants interact, it should rather observe what is
happening, understand why things work where they do work, why they do not work
where they do not work. Data analysis, modeling is a way of doing that.

Take, for instance, scaffolding. Scaffolding has a simple structure, made of
poles erected vertically, and horizontal boards attached to the poles that provide a
working surface for those using the scaffold. Scaffolding is simple to
understand, because we know how to build it, fitting tubes and boards together.
Now, imagine a leaf, with vascular bundles coming out of the stem and given the
leaft its particular shape. Leaves, obviously, have structure. They are not a
completely random assembly of the molecules that form the leaf. The structure of
leaves, however, is more difficult to comprehend than that of scaffolding.
Indeed, leaves are the result of many biological processes that lead to their
particular structure. And understanding why a leave has a particular shape
requires a lot of work on trying to uncover the relevant processes. 
Now, take cities. Cities are much more complicated than leaves.


Vascular systems are built once and for all and, if anything, only expand with
the body's size changing. On the other hand, cities adapt to the traffic on the
road which conversely is the result of the city's particular configuration.
Road evolve, buildings are built and destroyed, activities and people relocate,
etc. Cities are the delicate result of this interplay between transportation
networks and the activities, and a careful study of the former might lead to
invaluable insights into the structure and evolution of cities.




\subsection{We need data}
\label{sub:we_need_data}

Walk a few steps in your favourite city, feel the streets bustling all around
you. The sound of the cars, of people chatting, the pavement lined with
 homogeneously diverse buildings. The sense of familiarity we feel when stepping
back in a city that was once our home, years later. And that smell you had
forgotten you knew. Maybe the hardest thing, when studying cities, is the
impression that we know them closely. The belief that our impression of what
they are, the way we experience them, gives a true picture of what they really
are, the purpose they serve. In a sense, this is what makes the study of
macroscopic, human-made systems so difficult compared to the study of natural
systems. 

There are indeed only so many ways we can get acquainted with, say, electrons, and
therefore just so many things we can say about them. This, in a sense, makes the
study of electrons easy. Think about cities now. All the memories, habits,
knowledge you have gathered over the years. As individuals, we know too many and
too little things about them at the same time. We can have a very detailed
recollection of the city we have experienced. But this information is not
organised, and it is too local, too provincial. Therefore, we cannot infer what
cities are solely from our own experience.  We are a single piece of a puzzle
that counts hundreds of thousands, millions of them, all with a different
opinion of what their environment is like. No, to understand cities, how they
work as a system, we need to be told these thousands of stories, we need to
analyse them and see how similar, or dissimilar they really are. To understand
cities, we need data.\\


\subsection{Layers and scales}
\label{sub:layers_and_scales}

A simple and enlightening way to comprehend the extent of cities's complexity is
to consider the different layers and scales involved in the dynamics within and
of cities. At the smallest scale, we have the individuals who live in urban
systems. They make decisions about where they live, where they work, etc. and
interact constantly with one another. Individuals are, in a way, the building
blocks of cities, and it is therefore crucial to understand the way they
interact with their environment to understand the structure and behaviour of
cities.

At a larger scale, cities can be considered as systems characterised by specific
behaviours~\cite{Bettencourt:2007}. Besides, they do not evolve in isolation and
belong to larger scale structures. To quote the geographer B.J.L. Berry, `cities
[are] systems within systems of cities'~\cite{Berry:1964}, and their
interactions---migrations, commodity and capital flows---ought to constrain
their evolution~\cite{Pumain:2010}. 

Finally, there is a great amount of evidence to show that systems of cities also
exhibit very particular behaviours: the rank-size plot of the population of
cities that belong to these systems is indeed strikingly regular (a regularity
known as `Zipf's law'), and breaks down for other geographical units or when the
chosen set of cities is not geographically and economically
coherent~\cite{Cristelli:2012}.\\


\begin{figure}[!h]
    \centering
    \includegraphics[width=\textwidth]{./gfx/chapter-intro/spatial_scales.pdf}
    \caption{{\bf Interactions at different spatial scales.} Cities are the
    result of interactions occuring at different spatial scales. The movement
and interactions of individuals result in the properties of the city as a whole.
But cities are not closed systems, and interact with other cities in a system of
cities.\label{fig:spatialscale}}
\end{figure}


Cities are therefore the result of interactions occuring at different spatial
scales. Furthemore, they are not static: they evolve in time, through various
processes taking place at different time scales.

First we have time-scales of the order of a day, which span the daily commuting
of inhabitants. This incessant movement of people has been traditionally
explored through surveys, but new data now allow more thorough studies. The
digital traces that are left by people at all times (through their mobile phone,
metro pass or GPS device) indeed allow us to explore the structure of flows and
the pace of life in cities at unprecedently fine spatial and time resolutions.

Then, at the order of a year one can see the variation in terms of wealth,
population, etc. of cities, as recorded by statistics agencies. Data about
demographic, social and economic aspects of urban systems allow us to
characterise more specifically the structure and behaviours of these systems.

Finally, at time scales of the order of ten years, we can see the city's
infrastructure as well as its spatial footprint evolve. The study of the
underlying processes is made possible by various projects lead by the GIS
community, historians and geographers which aim at digitizing historical maps of
the road and rail networks in different regions of the world. Also, since the
$1970s$, many satellites have been taking pictures of the Earth's surface, and
the remote sensing community has been treating these data to get information
about the spatial extension of cities. These data should give us some insight
about the processes responsible for the long-term evolution of cities's
structure.\\

\begin{figure}[!h]
    \centering
    \includegraphics[width=0.9\textwidth]{./gfx/chapter-intro/time_scales.pdf}
    \caption{{\bf Different time scales.} The various data available about
    cities are associated with different time scales.\label{fig:timescale}}
\end{figure}

These time scales are summarised on Fig.~\ref{fig:timescale}. The long-term goal of our studies is to understand exactly how cities and systems of cities behave, and how interactions between these three layers lead to the behaviours we observe. 





% Different temporal and spatial scales at which they occur
A first step in the identification of order, is to identify the different
spatial and temporal scales at which phenomena are occuring. The goal of any
theory of how cities work would be to understand the phenomena occuring at each
scale, and link the phenomena occuring at different scales. Establish a
hierarchy of mechanism, as is the case in natural sciences~\cite{Simon:1962}.

%Dificulties with social sciences
How can we explain, then, the difference between social sciences and natural
ciences? Say, between economics and physics. Why, if the reductionist approach
is also to function with social systems, don't we have anything nearly as trong
a result as in the natural sciences? Clearly not because people are less
intellectually capable of handling difficult problems. Quite the
contrary, actually. One thing that jumps to your face when overlooking tebooks
in physics and textbooks in Geography, or in economics, is the difficulty of the
questions asked. The natural world is a complex realm as well, and there is no
reason to believe that the social world should be any more complicated (even
free minds obey to the law of large numbers when put together! And who tells you
electrons do not have a free mind too?). But the questions physics asks are
ridiculously simple, in comparison to what social scientists ask. In a way,
physicists have this natural advantage that Nature is a somewhat less immediate
reality to us, it is easier to distanciate oneself from physical phenomena than
it is from social phenomena. The problem with social sciences (if we can call
this a problem), is to try to solve complex questions because they are
\emph{important to us}. To try to develop an analogy, imagine physicists in the
18th century (before Kepler, before Newton...) were visited by an alien
civilisation, who would have described to them all the merits of the laser,
showed them one, and explained in what ways it could help us solve some basic
problems. Imagine they left, without leaving any explanation as to how build
one, or of the mecanisms involved behind the laser doing what it is doing.
Imagine politicians, pushed by the general public, decided to spend a lot of
money trying to build a laser. In what state do you think physics would be right
now? I am ready to bet no one would have never managed to build a laser, but we
we would have competing theories of what very general processes are at stake in
the processes of building a laser. Different schools of thoughts would have
emerged, urging to assemble this and this together. Until one day, maybe,
someone would have realised that we do not understand anything and have to go
back from basics, and forget about the laser for a time.
You only ever reach the goals you set to yourself.
Clearly, in some fields, economics being one of them, sociolog closely following
despite Bourdieu's criticism of the link between politics and sociology, this is
exactly what happening. Because they are asked to `make more money', `save the
conomy' or to justify such and such policy, these sciences have been preoccupied
with very high level, very complex questions that could be subdivided into tens,
hundreds, even more, simpler, more specific questions. They were asked to build
a rocket too soon. Successful approaches (with a few exceptions, see
thermodynamics) in physics were bottom-up approaches: by knowing the mecanisms
behind different phenomena, we were able to influence them, and combine them to
make objects that are more and more intricate. Top-down approaches rarely work,
because our brains are somehow not wired not handle very complex questions. But
we are cursed in the field of social sciences, because we are aware of the
top-down kind of questions, while you cannot aim at inventing a rocket before
knowing the basics of gravity. At least not seriously. There is something very
paradoxical in today's society in the very simple fact that while very complex
questions (regulate the economy, bringing full employment...) get massive
attention from funding agencies, while science-fiction-like projects such as
`inventing teletransportation' do not get a single euro from government
agencies. For me, they are of the same order. A wish to go towards an ideal,
allegedly better state of humanity, while not even knowing the mecanims on which
we would need to play in order to get there. While not even knowing
theoretically how to handle the problem. We only fund hard science projects that
we know are feasible, that are just a matter of technicalities, at least in
principle. Someone who would come up with the vague idea of writing a theory of
everything without any serious basis to think this is what he is doing would
have no chance to get any funding.\\


\begin{figure}
    \centering
    \includegraphics[width=1\textwidth]{gfx/chapter-intro/intra-urban.pdf}
    \caption{{\bf Intra-urban organisation.} Cities are first and foremost defined by
    the concentration of populations and various activities. The fact that
residences and activities have different locations is responsible for the
existence of flows of people, goods, etc. across the urban space. These flows
occur on appropriate infrastructure.\label{fig:intr_urban}}
\end{figure}

% Ackowledge the past
Of course, the description I have given of cities over the last few pages will
sound obvious to the reader owing to the work of predecessors.
